\documentclass[11pt]{article}
\usepackage[margin=1in]{geometry}
\usepackage{amsmath,amssymb,amsthm}
\usepackage[colorlinks=true,urlcolor=blue,linkcolor=blue]{hyperref}
\newtheorem{theorem}{Theorem}
\newtheorem{lemma}[theorem]{Lemma}
\newtheorem{proposition}[theorem]{Proposition}
\newcommand{\R}{\mathbb R}
\newcommand{\norm}[1]{\left\lVert#1\right\rVert}
\newcommand{\MM}{T_{\mathrm{MM}}}
\newcommand{\tO}{\widetilde O}
\newcommand{\diag}{\operatorname{diag}}
\newcommand{\range}{\operatorname{range}}
\newcommand{\rank}{\operatorname{rank}}
\title{Parallel product ridge actions: shared exact setup, fast blocks, and the unresolved filter step}
\author{Research working audit}
\date{November 14, 2026}
\begin{document}
\maketitle

\paragraph{Status.}
This note proves an exact-arithmetic improvement to the accounting of the
ridge primitive in \texttt{ridge-filters.tex}: all shifts share one
\(\tO(m^{\omega+\zeta})\) small-core setup, and true-system
corrections can use rectangular matrix multiplication on the entire block.
It does not prove the desired gap-free, high-accuracy rank-\(r\)
approximation of \(AB\). A large number of parallel shifts is affordable
in the supplied hard regime. The unresolved issue is stable evaluation or
subspace extraction at the accuracy required by a sharp filter. The algebraic
tools used here are established tools; novelty of their combination is not
certified.

\section{An exact shared-shift setup without an eigenvalue oracle}
Let \(P_0,Q_0\in\R^{m\times m}\) be positive semidefinite. They
need not be nonsingular. For \(\lambda>0\), use the unscaled block
system
\begin{equation}\label{eq:unscaled}
 U_\lambda=
 \begin{pmatrix}P_0&I\\-\lambda I&Q_0\end{pmatrix}.
\end{equation}
This is a rescaling of the block system in \texttt{ridge-filters.tex}:
with \(w=\sqrt\lambda\,y\), the two formulations are equivalent.
The unscaled form keeps the new setup and applications entirely within
field operations \(+, -,\times,/\).

\begin{lemma}[Shared exact resolvents]\label{lem:shared}
Given \(P_0,Q_0\), a deterministic setup costing
\(\tO(m^{\omega+\zeta})\) field operations permits the following
operation. For any positive shifts \(\lambda_1,\ldots,\lambda_s\)
and arbitrary right-hand sides \((h_j,k_j)\), all solutions
\[
 U_{\lambda_j}\binom{x_j}{w_j}=\binom{h_j}{k_j}
\]
can be computed exactly using a constant number of
\(\MM(m,m,s)\) products and \(O(ms)\) additional field operations.
The shifts and right-hand sides may be chosen after setup.
\end{lemma}

\begin{proof}
Form \(D=P_0Q_0\). Neiger--Pernet--Villard \cite{npv}, Corollary 5.2,
compute a Frobenius normal form together with its transformation
deterministically in
\(O(m^\omega(\log\log m)^2)\) field operations. Thus, after also
inverting the transformation, we have
\[
 D=S F S^{-1},
\]
where \(F\) is block diagonal with companion blocks. This is an exact
field computation. It does not require computing generally irrational
eigenvalues by finitely many rational operations.

For a companion block of dimension \(d\), solving
\((F_i+\lambda I)u=v\) costs \(O(d)\) operations. For example,
take the convention with ones on the subdiagonal and the negated monic
polynomial coefficients in the last column. The last \(d-1\) row
equations express \(u_{d-1},\ldots,u_1\) affinely in \(u_d\),
in backward order. The first row then determines \(u_d\). The
denominator is nonzero precisely when the shifted companion matrix is
nonsingular. This is a Horner-type recurrence; it does not require a
dense triangular solve. Summing over blocks gives \(O(m)\) per
right-hand side.

For \(\lambda>0\), \(D+\lambda I\) is nonsingular. Indeed,
the nonzero eigenvalues of \(P_0Q_0\) are nonnegative, as seen by
comparison with \(Q_0^{1/2}P_0Q_0^{1/2}\). Zero eigenvalues are
also harmless. Consequently all companion recurrences are valid.

Eliminate \(x\) from \eqref{eq:unscaled}. The result is
\begin{equation}\label{eq:block-solution}
 w=(D+\lambda I)^{-1}(\lambda h+P_0k),\qquad
 x=\lambda^{-1}(Q_0w-k).
\end{equation}
For all columns, first form \(H\diag(\lambda_j)+P_0K\), change
basis by \(S^{-1}\), solve the companion blocks independently at each
shift, change basis by \(S\), and multiply by \(Q_0\). These are
a constant number of common dense matrix--block products and \(O(ms)\)
scalar operations. The setup is independent of the shifts.
\end{proof}

This observation corrects the unnecessary per-shift setup charge in the
earlier corollary. It is an improvement in the proved primitive's accounting,
not a claim that simultaneous shifted solves or Frobenius forms are new.
The algebraic setup is deterministic, so it introduces no new failure event.

\paragraph{Equivalent spectral formula, including singular Gram blocks.}
For comparison, suppose spectral data are available and write
\[
 Q_0^{1/2}P_0Q_0^{1/2}=W\diag(d_i)W^\top,
 \quad G=Q_0^{1/2}W,\quad F_0=P_0G,\quad
 D_\lambda=\diag((d_i+\lambda)^{-1}).
\]
For the original scaled block matrix
\(T_\lambda=\bigl(\begin{smallmatrix}P_0&\sqrt\lambda I\\
-\sqrt\lambda I&Q_0\end{smallmatrix}\bigr)\), direct substitution gives
\[
 T_\lambda^{-1}=
 \begin{pmatrix}
 GD_\lambda G^\top&
 -\lambda^{-1/2}(I-GD_\lambda F_0^\top)\\
 \lambda^{-1/2}(I-F_0D_\lambda G^\top)&
 \lambda^{-1}(P_0-F_0D_\lambda F_0^\top)
 \end{pmatrix}.
\]
No inverse of \(P_0\) or \(Q_0\) occurs. Shared basis transforms
and per-column Cauchy denominators likewise give a common setup for all
shifts. But an exact eigendecomposition is not a finite exact-field
primitive; Lemma~\ref{lem:shared} supplies the strictly algebraic route.

\section{Fast rectangular corrections and a uniform error theorem}
Let \(A\in\R^{n\times m}\), \(B\in\R^{m\times p}\),
\(P=A^\top A\), \(Q=BB^\top\), \(I=m(n+p)\), and
\(M=AB\). Suppose constant-distortion sketches satisfy
\[
 (1-\eta)P\preceq\widetilde P\preceq(1+\eta)P,
 \qquad
 (1-\eta)Q\preceq\widetilde Q\preceq(1+\eta)Q,
 \quad 0<\eta<1/2.
\]
Define, up to a fixed multiplicative constant, the batched application cost
\begin{align*}
 \mathcal B(s)={}&\MM(n,m,s)+\MM(m,n,s)\\
 &+\MM(p,m,s)+\MM(m,p,s)+\MM(m,m,s).
\end{align*}
The true Gram actions on a block are computed as
\(PX=A^\top(AX)\) and \(QY=B(B^\top Y)\); no true full Gram
matrix is formed.

\begin{theorem}[Uniform batched ridge actions]\label{thm:batch}
After \(\tO(I+m^{\omega+\zeta})\) setup, a block of \(s\)
possibly different positive-shift ridge actions
\[
 R_{\lambda_j}v_j=
 MM^\top(MM^\top+\lambda_j I)^{-1}v_j
\]
can be approximated to
\[
 \norm{\widehat R_{\lambda_j}v_j-R_{\lambda_j}v_j}_2
 \le\xi\norm{v_j}_2\qquad(j=1,\ldots,s)
\]
in \(O(\mathcal B(s)\log(1/\xi))\) additional arithmetic
operations. Logarithmic constants depend only on \(\eta\).
The embedding event makes this statement simultaneous for all shifts,
right-hand sides, iteration counts, and adaptive choices.
\end{theorem}

\begin{proof}
Use Lemma~\ref{lem:shared} with
\(P_0=\widetilde P\), \(Q_0=\widetilde Q\) to apply every
sketched block inverse after one shared setup. Apply the Richardson
correction from the uniform ridge theorem in \texttt{ridge-filters.tex}
to all columns in parallel; the rescaling between \(U_\lambda\)
and \(T_\lambda\) preserves the exact iterates after the same
coordinate change. The proven error estimate is
\[
 \norm{\widehat R_{\lambda,j}-R_\lambda}_2
 \le K_\eta\rho_\eta^j,\qquad
 K_\eta=\sqrt{\frac{1+\eta}{1-\eta}},\quad
 \rho_\eta=\frac\eta{1-\eta}<1.
\]
Each simultaneous true-system action costs \(O(\mathcal B(s))\).
Each simultaneous preconditioner action costs
\(O(\MM(m,m,s)+ms)\), which is included in the same bound.
Initial products \(A^\top[v_1,\ldots,v_s]\) and final lifting
by \(A\) are also included. Choose
\(j=O(\log(K_\eta/\xi))\).
\end{proof}

The input sketches can each be assigned failure probability \(\delta/2\).
Once both inequalities hold, the theorem is deterministic and simultaneous;
there is no union bound over the \(s\) shifts or repeated adaptive blocks.
Rank deficiency adds no new parameter, by the rank-deficient ridge theorem
and Lemma~\ref{lem:shared}.

For \(L\) shifts each applied to \(r\) vectors, use \(s=Lr\)
columns. The new bound is
\[
 \tO(I+m^{\omega+\zeta})+
 O\!\left(\mathcal B(Lr)\log(1/\xi)\right),
\]
instead of requiring \(L\) independent fast factorizations. Memory for
the returned block and any subsequent low-rank factors must still be paid.
This is a vector-action theorem; it does not output the AMM approximation.

\section{Hard-regime block accounting}
Take \(n=p=N\), \(m=N^{8/9}\), so that \(N=m^{9/8}\).
A block of width \(s=m^{1/2}\) has matrix-multiplication dimension
exponents
\[
 (a,b,c)=(9/8,1,1/2).
\]
The full rectangular bound in Amanov--van den Berg--Zuiddam
\cite{avz}, Theorem 1.2, is input-optimal when
\(c(a+b)\le ab\). Here
\[
 c(a+b)=17/16\le9/8=ab.
\]
Thus
\[
 \mathcal B(\sqrt m)=m^{17/8+o(1)}=N^{17/9+o(1)}.
\]
In particular \(\sqrt m\) simultaneous exact Gram actions can be
computed at the input-size exponent. The shared setup costs
\(m^{9/4+o(1)}=N^{2+o(1)}\). Polylogarithmically many refinement
rounds with polynomially small \(\xi\) therefore fit the desired
\(N^{2+o(1)}\) construction exponent. The number of positive shifts
alone is not an obstruction.

More generally, a block of width \(m^c\) remains input-optimal for
\(c\le9/17\). For larger \(c\), the same theorem gives
\[
 \omega(9/8,1,c)\le\frac{(9+17c)^2}{288c}.
\]
This is at most \(9/4\) through
\(c\le(171+54\sqrt2)/289\), approximately \(0.856\).
This larger block may fit the total small-core budget even though it is
not input-optimal. These statements concern one known block of right-hand
sides, not sequentially dependent queries that have yet to be generated.

If an algorithm eventually returns width-\(r\) factors, it must add
\(O(npr)\) for explicit expansion. In the hard rank-one regime that
adds \(O(N^2)\). None of the batched-action statements eliminates this
output cost.

\section{Why residual augmentation does not yet repair a sharp filter}
For a signed linear combination of negative-shift ridge actions,
\(f(H)=\sum_\ell a_\ell R_{\lambda_\ell}\), let
\(\mathcal A=\sum_\ell|a_\ell|\). The usual separate-action bound
is
\[
 \norm{\widehat f(H)-f(H)}_2\le\mathcal A\xi.
\]
The coefficient theorem in \texttt{ridge-filters.tex} permits
\(\log\mathcal A=\Omega(\theta^{-1/2}\log m)\) for a sharp
uniform cutoff with polynomially small tail leakage. This theorem does
not prove a lower bound against another evaluation method. It does imply
that Theorem~\ref{thm:batch}, used with this separate-action error bound,
needs polynomially many refinement rounds in the hard example.

If \(y_\ell\) approximately solves
\((H+\lambda_\ell I)y_\ell=b\), write
\[
 r_\ell=b-(H+\lambda_\ell I)y_\ell.
\]
Adjoining the residuals to the candidate space gives the exact relation
\[
 H y_\ell\in\operatorname{span}\{b,y_\ell,r_\ell\}.
\]
This is useful algebra, but it does not identify that enlarged space with
the exact high-degree rational Krylov space. The right-hand sides have
changed independently at each shift. A theorem controlling the amplified
combination, or an alternative rank-\(r\) extraction theorem, is still
needed.

\begin{proposition}[Scalar-preconditioner test]\label{prop:scalar}
Consider simultaneous Richardson solves of
\((H+\lambda I)y=b\), initialized at zero, with preconditioner
\((c+\lambda)I\), for any number of different positive shifts.
After at most \(j\) correction rounds, all iterates and their
residuals lie in
\[
 \mathcal K_{j+1}(H,b)
 =\operatorname{span}\{b,Hb,\ldots,H^j b\}.
\]
Increasing the number of shifts does not increase this upper bound on
the exact generated space.
\end{proposition}
\begin{proof}
Every correction is a scalar multiple of a previous residual. Induction
shows that an iterate after \(j\) rounds is a polynomial in \(H\)
of degree at most \(j-1\), applied to \(b\); its residual has degree
at most \(j\). Different shifts change only the scalar coefficients.
\end{proof}

Take \(H\) with spectrum in a constant relative interval around
\(c>0\). Then \(cI\) is an admissible constant-distortion
preconditioner, so Proposition~\ref{prop:scalar} is a valid test for any
argument using only that distortion property. The same issue occurs in
the product block formulation with \(Q=\widetilde Q=I\) and
\(\widetilde P=cI\): the blocks of the sketched inverse are scalar
multiples of the identity, and all output vectors have polynomial degree
controlled by the number of correction rounds, not the number of shifts.
This is not an OSE distributional lower bound. A proof exploiting the
actual randomness of the input embedding could go beyond this test, but
must explicitly establish the required additional property.

Stable sequential rational Arnoldi or a factorized Newton/Clenshaw
evaluation avoids forming a huge partial-fraction sum directly. Its
known operator dependencies are sequential, so it does not automatically
inherit the fast-block budget. A simultaneous Sylvester or Clenshaw
linearization is an interesting alternative. A constant sketch can move
its surrogate eigenvalues outside the Chebyshev interval by a fixed
amount, however; high-degree polynomial values there can grow
exponentially in the degree. No uniform contraction or stable fast
parallel solve for that linearization has been proved here.

\section{Exact partial Gram postprocessing and the sequential gap}
Suppose a coefficient block \(V\in\R^{m\times s}\) is already
available, and let \(Z=AV\). Exact restricted information can be
obtained through
\[
 Z=AV,\quad F=A^\top Z,\quad T=B^\top F=M^\top Z,
 \quad G=Z^\top Z,\quad K=T^\top T.
\]
The two small Gram matrices supply normalization of \(Z\) and the
usual Rayleigh--Ritz core for the product restricted to its span. Only
the selected \(r\) directions need be lifted. The required products
cost a constant number of the terms in \(\mathcal B(s)\), plus
\(\MM(s,n,s)+\MM(s,p,s)\), fast \(s\)-dimensional algebra,
and selected-direction lifting. For \(s=\sqrt m\) in the hard
regime, all these costs fit \(N^{2+o(1)}\).

For a genuine standard Krylov candidate, its established low-rank
approximation proof justifies the final selection. Merely knowing a
good rank-\(s\) range approximation is not sufficient to assert the
same \((1+\theta)\) error after truncating to rank exactly \(r\);
the rank-\(r\) extraction needs its own guarantee.

The unresolved construction is the true coefficient Krylov block
\[
 V=[\Omega,\,QP\Omega,\,(QP)^2\Omega,\ldots].
\]
Its columns are not known before the exact input actions are performed.
Charging one fast multiplication to the entire final block therefore
does not constitute an algorithm for generating it. With conventional
sequential generation and \(s=\Theta(\theta^{-1/2})=\Theta(\sqrt m)\)
for rank one, the input work remains
\(\tO(I\sqrt m)=N^{7/3+o(1)}\), above both the target and the
named fast-Gram comparison \(N^{169/81+o(1)}\).

Neiger--Pernet--Villard's fast Krylov theorem applies to an
\emph{explicit} square operator. Applying it to \(QP\) after
forming \(P,Q\) retains the Gram setup cost. Treating the factorized
operator as if that explicit input had already been supplied would
discard the central cost in this problem.

\paragraph{A genuinely different remaining route.}
Their polynomial-matrix, kernel-basis, and high-order-lifting tools might
be applied directly to a tall factor linearization. For example, the
polynomial matrix
\[
 \mathcal L(z)=
 \begin{pmatrix}
 zI_m&-A^\top&0&0\\
 0&I_n&-A&0\\
 0&0&I_m&-B\\
 -B^\top&0&0&I_p
 \end{pmatrix}
\]
has Schur complement \(zI_m-PQ\) after eliminating the identity
blocks. Its polynomial degree is concentrated in only \(m\) rows,
while its off-diagonal data are precisely the tall factors. A
structure-sensitive kernel-basis or truncated-power-projection theorem
for this representation could be relevant. Standard explicit Schur
elimination forms the costly Gram blocks, and a dense algorithm on the
full \((n+p+2m)\)-dimensional linearization is also too expensive.
No appropriate cost theorem for retaining this structure is proved or
identified here. The explicit-square-matrix theorem must not be claimed
to cover it.

\section{Precision and final assessment}
The shared-shift theorem is a proved exact-field arithmetic result. It
avoids the logical error of charging an exact spectral decomposition as
a finite rational computation. It does not make the canonical-form
similarity basis well conditioned. Companion coefficients, transformation
matrices, and nearly cancelling terms in \eqref{eq:block-solution}
can be large; shifts near zero can require very high relative precision.
The uniform output-norm proof for exact Richardson does not by itself
control those internal roundoff errors. A numerical implementation needs
an explicit precision theorem, potentially depending on shift scale and
core conditioning. No polylogarithmic bit-cost result follows from the
operation count alone.

The positive conclusions are one shared exact setup for every shift and
an input-optimal exponent for a \(\sqrt m\)-wide known block in the
hard regime. The arbitrary-input AMM target remains unproved. There is
no demonstrated new total-runtime algorithm beating the named
\(N^{169/81+o(1)}\) comparison on every input. The most concrete new
question isolated here is whether structured factor linearizations or a
structure-preserving parallel rational-space extraction can exploit the
affordable batch without inheriting sequential generation or exponential
cancellation precision.

\begin{thebibliography}{9}
\bibitem{npv}
V. Neiger, C. Pernet, and G. Villard.
\emph{Computing Krylov Iterates in the Time of Matrix Multiplication}
(2024), Theorem 4.1 and Corollary 5.2.
\url{https://arxiv.org/abs/2402.07345}.
\bibitem{avz}
A. Amanov, M. van den Berg, and J. Zuiddam.
\emph{Multiplying Matrices via Polynomials: Exposition and Barriers}
(2026 preprint), Theorem 1.2.
\url{https://arxiv.org/abs/2610.30054}.
\bibitem{mm}
C. Musco and C. Musco.
\emph{Randomized Block Krylov Methods for Stronger and Faster Approximate
Singular Value Decomposition} (2015).
\url{https://arxiv.org/abs/1504.05477}.
\end{thebibliography}
\end{document}
